% SLICEDENSITY2
\documentclass[11pt]{article}
\usepackage[margin=.9in]{geometry}
\usepackage{amsmath,amssymb,amsthm}
\begin{document}
\begin{center}
{\Large SLICEDENSITY2: Move the primary slice above the pure-literal regime}\\[4pt]
\textbf{Date: 2026-10-04}\\
\textbf{Status: CANDIDATE roadmap correction}\\
\textbf{P versus NP status: OPEN}
\end{center}

\section{Correction to SLICEDENSITY1}

SLICEDENSITY1 correctly proves computational hardness for every constant exact-layer density
\[
c>2/3
\]
and correctly transfers the rigorous random-3SAT satisfiability statement for
\[
c<3.52
\]
to the distinct-clause exact layer.

However, its recommendation \(c=1\) is too aggressive.

Scott and Sorkin show that below the pure-literal threshold, rare UNSAT random formulas are typically explained by a unique minimal unsatisfiable subformula of the smallest possible excess, and they give efficient algorithms in sufficiently sparse regimes. Molloy's core analysis identifies the pure-literal transition; for random 3-SAT its value is about
\[
c_{PL}\approx1.636.
\]
Thus \(c=1\) lies in a regime where conditioned UNSAT instances are known to have local/easy structure. That is undesirable for a positive test distribution.

\section{Revised density window}

The useful rigorously controlled interval is
\[
\boxed{c_{PL}<c<3.52}.
\]
In this window:
\begin{itemize}
\item[] exact-layer SAT remains NP-complete and exact-layer UNSAT coNP-complete by the control-variable padding theorem;
\item[] the uniform exact layer is SAT with high probability by Hajiaghayi--Sorkin plus the negligible-collision transfer;
\item[] the Scott--Sorkin small-MUF characterization below the pure-literal threshold no longer applies;
\item[] the structure of rare UNSAT formulas between the pure-literal and satisfiability thresholds is explicitly identified in their work as unresolved.
\end{itemize}

A concrete working value is
\[
\boxed{c=3.4,\qquad t=\lceil3.4q\rceil}.
\]
This stays strictly below the rigorous \(3.52\) satisfiable regime while lying well above the pure-literal threshold and the classical simplest peeling regime.

\section{Why this is only a target selection}

No theorem currently says that rare UNSAT formulas at density \(3.4\) have the globally distributed, large-core structure needed for a monotone lower bound. Scott--Sorkin explicitly leave the intermediate region unresolved.

Therefore the next positive-distribution theorem is itself substantive:
\[
\boxed{
\text{INTERMEDIATE-CORE1: characterize rare UNSAT exact-}\lceil3.4q\rceil
\text{ formulas}
}
\]
with the specific objective of ruling out polynomially recognizable local obstructions and finding a large spread subfamily with matched local statistics.

\section{Interaction with DENSITYPOSFAIL1}

The P/POLY-POSITIVE-BARRIER still applies. A hand-picked positive family with polynomial-size membership circuits cannot support a distributional monotone lower bound.

The degree-profile hostile control also remains mandatory. Any planted positive model must match:
\begin{itemize}
\item[] variable degree distribution;
\item[] low-order clause intersection/codegree statistics;
\item[] small-core/excess statistics;
\item[] sign/polarity marginals.
\end{itemize}

\section{Ranked plan}

\begin{enumerate}
\item \textbf{INTERMEDIATE-CORE1:} analyze rare UNSAT formulas in the exact density \(c=3.4\) layer, computationally first at moderate \(q\), measuring minimal UNSAT core size, deficiency, variable degrees, codegrees, and incidence expansion.
\item \textbf{SAT-SUNFLOWER1:} only after identifying a positive hard-core distribution whose local statistics overlap the uniform SAT layer.
\item \textbf{SLICELIFT-NONLOCAL:} secondary route; SLICELIFT2 rules out block-local exact-weight gadgets.
\item \textbf{Restricted models:} continue only as theorem-building laboratories.
\end{enumerate}

\section{Status}

This note corrects the preferred density, not the hardness theorem. The exact-layer lower bound remains OPEN and \(P\ne NP\) remains OPEN.

\begin{thebibliography}{4}
\bibitem{ScottSorkin2010}
A. D. Scott and G. B. Sorkin, Structure of random r-SAT below the pure literal threshold, 2010.
\bibitem{Molloy2005}
M. Molloy, Cores in random hypergraphs and Boolean formulas, \emph{Random Structures and Algorithms} 27(1):124--135, 2005.
\bibitem{HajiaghayiSorkin2003}
M. T. Hajiaghayi and G. B. Sorkin, The Satisfiability Threshold of Random 3-SAT Is at Least 3.52, 2003.
\bibitem{Valiant1986}
L. G. Valiant, Negation is Powerless for Boolean Slice Functions, \emph{SIAM Journal on Computing} 15(2):531--535, 1986.
\end{thebibliography}
\end{document}
